\section{Methodology}
\label{section:algorithm}

Our methodology follows a two-stage framework built upon a pre-trained expressive IndexTTS2 model. Starting from a model trained on a high-resource language, we first adapt it to the low-resource target language, and then apply reinforcement learning to preserve and recover emotional expressiveness, prosodic control, and content accuracy during transfer.

\subsection{Stage 1: Low-Resource Language Adaptation}
In the first stage, we adapt the expressive base model to the low-resource target language. The goal of this stage is to acquire target-language pronunciation and linguistic patterns while preserving as much of the pretrained expressive prior as possible. To reduce severe catastrophic forgetting, we adopt a parameter-efficient adaptation strategy, such as lightweight adapter modules or LoRA, so that only a limited subset of model parameters is updated during target-language transfer.

\subsection{ASQP Annotation Pipeline}
\label{section:algorithm:backtranslation}
The Valence--Arousal (VA) conditioning and local emphasis annotations required for the \emph{VA-and-stress control} are derived from an Aspect Sentiment Quad Prediction (ASQP) model, which operates on Traditional Chinese text. The RL training data is sourced directly from the SemEval-2026 Task 3 Chinese restaurant dataset~\citep{DimABSA2026}, which provides Traditional Chinese sentences with aspect--sentiment quadruple annotations and thus requires no translation.

For the evaluation corpus (TAT-Vol1 Taiwanese test utterances), however, no Chinese-language annotations are available. Since available ASQP tools do not support the romanized Taiwanese representation, we adopt a back-translation pipeline for this portion only: each Taiwanese test sentence is first translated into Traditional Chinese, processed by the ASQP model to extract aspect--sentiment quadruples, and then translated back to Taiwanese. The reverse translation step employs the heterogeneous model ensemble with adaptive weighted voting proposed in our previous work~\citep{Hsieh2026NYCU}, which aggregates multiple translation candidates to improve lexical fidelity under the low-resource constraint. The resulting sentences, together with their ASQP-derived annotations, form the \textit{Back-Translated Dataset} used to evaluate IndexTTS2-FT-va-emphasis, IndexTTS2-FT+RL, and IndexTTS2-proposed.

\subsection{Phonemic Text Front-End}
Because the low-resource target language is written with Han characters but pronounced according to its own phonology, we drive the model with a romanized phonemic representation rather than raw Han text. Each input sentence is converted to Taiwanese romanization (T\^ai-l\^o) with numeric tone marks (e.g., \texttt{gua2 huann1-hi2}), and locally emphasized spans are wrapped in quotation markers that the model learns to render as prosodic stress, providing the textual signal that the Local Emphasis Control Reward acts on. Crucially, the backbone's built-in text normalizer---designed for Mandarin/English number and symbol expansion---is bypassed, since the input is already phonemic and normalization would corrupt both the romanization (including the tone digits) and the emphasis markers. The identical front-end is applied during RL training and inference, so the representation the policy is optimized on exactly matches the one used at generation time.

\subsection{Stage 2: RL-Based Emotional Realignment}
In the second stage, we introduce an RL-based alignment procedure~\citep{2025arXiv251005758L} on top of the adapted model. The core idea is to treat the autoregressive generation of semantic speech representations as a sequential decision-making process, while computing task-specific rewards from the decoded speech waveform. This stage explicitly optimizes emotion preservation, emotional intensity, local emphasis, and content intelligibility in the low-resource language. In practice, the RL updates are mainly applied to the autoregressive T2S transformer, and can optionally include the emotion- or style-related conditioning modules when stronger expressive realignment is needed.

\begin{figure}[htbp]
    \centering
    \includegraphics[width=0.8\linewidth]{asset/architecture.pdf}
    \caption{High-level architecture of the proposed RL-enhanced IndexTTS2 framework. After low-resource adaptation, the generator outputs semantic speech tokens, which are decoded into speech waveforms and evaluated by emotion, intensity, prosody, and content reward models. These rewards are used to update the policy via Group Relative Policy Optimization (GRPO), mitigating emotional forgetting during language transfer.}
    \label{fig:architecture}
\end{figure}

\subsection{Reinforcement Learning Formulation}
The generation of semantic speech tokens is formulated as a Markov Decision Process (MDP), while the reward is computed from the decoded speech waveform:
\begin{itemize}
    \item \textbf{State ($S_t$):} The hidden states of the IndexTTS2 autoregressive decoder, incorporating the input text, speaker conditions, and previously generated speech tokens up to step $t$.
    \item \textbf{Action ($A_t$):} The predicted semantic speech token at step $t$, sampled from the policy distribution $\pi_\theta(A_t|S_t)$.
    \item \textbf{Reward ($R_t$):} A composite reward signal provided upon the completion of a sequence or segment. The reward is designed to preserve emotional expressiveness while maintaining intelligibility in the low-resource target language, and comprises four components:
    \begin{enumerate}
        \item \textit{Emotion Classification Reward:} Evaluated using a pre-trained emotion recognition model (e.g., emotion2vec) to measure whether the generated speech matches the target emotion category.
        \item \textit{Global Emotion Intensity Reward:} Computed from the distance between the generated speech and a neutral centroid in the Valence-Arousal-Dominance (VAD) space, in order to encourage the target global emotional intensity.
        \item \textit{Local Emphasis Control Reward:} A reward based on the relative pitch and energy patterns of emphasized words, which encourages correct placement of local stress and expressive prosody.
        \item \textit{ASR-Based Content Reward:} Computed using an ASR model on the generated speech to compare the transcription against the input text, encouraging content preservation and intelligibility during RL optimization.
    \end{enumerate}
\end{itemize}
In summary, the sequence-level reward can be written as
\[
R = \lambda_1 R_{\text{emo-cls}} + \lambda_2 R_{\text{intensity}} + \lambda_3 R_{\text{emphasis}} + \lambda_4 R_{\text{ASR}},
\]
where the coefficients $\lambda_i$ control the trade-off between emotional expressiveness, local prosody, and content accuracy.

\subsection{Alleviating Catastrophic Forgetting}
To address the degradation of emotional capabilities during low-resource fine-tuning, we initialize our policy $\pi_\theta$ with the weights of the target-language adapted IndexTTS2 model rather than training from scratch. We then employ Group Relative Policy Optimization (GRPO) mixed with a Kullback-Leibler (KL) divergence penalty against the adapted reference model. This ensures that the model explores localized emotional and prosodic improvements while preserving the linguistic adaptation and expressive prior inherited from the original high-resource model.
